\chapter{Introducción}
\label{cap:introduccion}

\chapterquote{Many things we can't do are simply because we think we can't do them}{Satoru Iwata}

El estudio de grandes corpus musicales constituye una de las principales herramientas de investigación dentro de disciplinas como la musicología histórica, la etnomusicología, el análisis musical o la documentación del patrimonio sonoro. A través del análisis sistemático de conjuntos extensos de obras es posible identificar patrones estilísticos, estructuras compositivas recurrentes, relaciones entre repertorios y fenómenos culturales que difícilmente pueden observarse mediante el estudio aislado de piezas individuales.

Tradicionalmente, este tipo de investigaciones ha dependido de procesos manuales de lectura, catalogación y análisis de partituras, una metodología rigurosa pero limitada en términos de escalabilidad. Asimismo, este enfoque analógico (e incluso algunas aproximaciones digitales estructuradas convencionales) introduce una importante barrera de rigidez metodológica; cuando surgen nuevas preguntas de investigación o hipótesis imprevistas, se genera la necesidad de volver a catalogar el corpus por completo y de reestructurar la arquitectura de las bases de datos pertinentes para incluir las nuevas variables. El creciente volumen de material musical digitalizado, especialmente en ámbitos relacionados con la música tradicional y el patrimonio musical, ha puesto de manifiesto la necesidad de desarrollar herramientas capaces de automatizar parte de estas tareas, flexibilizar el uso de la información y facilitar la exploración dinámica de grandes colecciones documentales.

La aparición de formatos simbólicos estandarizados para la representación musical, como MusicXML \citep{good2001musicxml} o MEI \citep{roland2002music} ha impulsado el desarrollo de la musicología computacional y ha permitido aplicar técnicas de análisis automatizado sobre repertorios cada vez más amplios. Sin embargo, la práctica demuestra que siguen existiendo importantes desafíos relacionados con la heterogeneidad de los corpus, la compatibilidad entre formatos, la extracción de información musicológica compleja, la integración entre contenido musical y textual, y la recuperación eficiente de conocimiento a partir de grandes volúmenes de datos.

De forma paralela, la aparición de los LLMs y los avances recientes en el procesamiento del lenguaje natural han abierto un nuevo paradigma en la interacción con colecciones de datos extensas. Al contrario que los sistemas indexados tradicionales, los LLMs ofrecen el potencial de suavizar la rigidez metodológica de las bases de datos convencionales, permitiendo formular preguntas de investigación imprevistas en lenguaje natural sin la necesidad de reestructurar manualmente los esquemas relacionales o de volver a catalogar exhaustivamente el corpus desde cero. No obstante, las particularidades de la información musical simbólica plantean limitaciones severas que dificultan la aplicación directa de estos modelos. Debido a la densidad de formatos como \citep{good2001musicxml} o MEI \citep{roland2002music}, la inyección directa de partituras satura con rapidez las ventanas de contexto de los modelos. Además, al carecer de capacidades nativas para el cálculo estructurado, los LLMs resultan ineficaces a la hora de resolver consultas que requieran conocimientos derivados de análisis musicales especializados (tales como la detección de síncopas, hemiolas o motivos musicales) que no se encuentran explícitamente representados en las fuentes textuales originales, lo que tiende a generar alucinaciones factuales y restringe su utilidad en la investigación musicológica técnica si no se integran en arquitecturas híbridas y especializadas.

En este contexto, este trabajo presenta el diseño, desarrollo e implementación de una arquitectura híbrida de ingeniería basada en inteligencia artificial orientada al análisis automatizado y la consulta flexible de corpus musicales folclóricos extensos y heterogéneos. Para superar el sesgo de las tecnologías generativas puras, el sistema despliega un pipeline de ingesta capaz de procesar colecciones simbólicas (\textit{MusicXML}, \textit{MEI} y \textit{MSCZ}), extrayendo y normalizando mediante la librería \textit{music21} un espectro de variables de alta complejidad que abarca desde metadatos básicos hasta fenómenos que requieren un análisis musical más complejo como síncopas, hemiolias horizontales y verticales, motivos musicales y control de calidad estructural. La originalidad de la solución propuesta radica en el diseño y desarrollo de una infraestructura local determinista y libre de alucinaciones factuales fundamentada en una base de datos relacional \textit{SQLite} interconectada a un modelo de lenguaje pequeño (\textit{Llama3.1:8B}) mediante el estándar avanzado \textit{Model Context Protocol} (MCP) \citep{anthropic2024mcp} para la traducción automática de lenguaje natural a consultas estructuradas (\textit{Text-to-SQL}). Con el fin de evaluar el rendimiento y validar la eficacia de esta propuesta, la investigación incluye un estudio comparativo riguroso frente a una implementación basada en RAG mediante el motor nativo \textit{File Search} de \textit{Gemini}, demostrando empíricamente cómo el enfoque local estructurado por MCP supera las limitaciones semánticas y las tasas de alucinación de los modelos comerciales de gran escala ante consultas musicológicas complejas. Esta solución se expone de forma accesible a través de una interfaz de aplicación web, dotando a la comunidad etnomusicológica de una herramienta para la exploración dinámica y la extracción de información en grandes colecciones musicales.

Para dar respuesta a estos desafíos de flexibilidad analítica, automatización de procesamiento de formatos simbólicos y superación de las limitaciones de la IA generativa, esta arquitectura se enmarca y desarrolla dentro del proyecto europeo e iberoamericano \textit{``EA-DIGIFOLK — An European and Ibero-American approach for the digital collection, analysis and dissemination of folk music (101086338)''}, financiado por la Unión Europea y en el cual participa activamente la Universidad Complutense de Madrid (UCM)\footnote{\url{https://cordis.europa.eu/project/id/101086338}}. Si bien el propósito de \textit{EA-DIGIFOLK} es preservar, compartir y difundir el patrimonio folclórico musical popular en entornos públicos y educativos a través de una plataforma accesible, el núcleo metodológico fundamental que sustenta dicho objetivo reside en el análisis musical y etnomusicológico profundo del material audiovisual y escrito recopilado. Es en este aspecto científico especializado donde este trabajo de investigación cobra sentido, proporcionando la infraestructura tecnológica necesaria para transformar el corpus digitalizado en conocimiento explícito, auditable y dinámico.

% El estudiante elaborará una memoria descriptiva del trabajo realizado, con una \textbf{extensión mínima recomendada de 50 páginas} incluyendo al menos una introducción, objetivos y plan de trabajo, resultados con una discusión crítica y razonada de los mismos, conclusiones y bibliografía empleada en la elaboración de la memoria.

% La memoria se puede redactar en castellano o en inglés, pero en el primer caso la introducción y las conclusiones de la memoria tienen que traducirse también al inglés y aparecerán como capítulos \textbf{al final de la memoria}. En ambos casos, el título de la memoria aparecerá en castellano y en inglés.

% Además del cuerpo principal describiendo el trabajo realizado, la memoria contendrá los siguientes elementos, que no computarán para el cálculo de la extensión mínima del trabajo:

% \begin{itemize}
% \item un resumen en inglés de media página, incluyendo el título en inglés,
% \item ese mismo resumen en castellano, incluyendo el título en castellano,
% \item una lista de no más de 10 palabras clave en inglés,
% \item esa misma lista en castellano,
% \item un índice de contenidos, y
% \item una bibliografía.
% \end{itemize}


% Para facilitar la escritura de la memoria siguiendo esta estructura, el estudiante podrá usar las plantillas en LaTeX o Word preparadas al efecto y publicadas en la página web del máster correspondiente.

% Todo el material no original, ya sea texto o figuras, deberá ser convenientemente citado y referenciado. En el caso de material complementario se deben respetar las licencias y copyrights asociados al software y hardware que se emplee. En caso contrario no se autorizará la defensa, sin menoscabo de otras acciones que correspondan.

\section{Motivación}
% \subsection{Limitaciones actuales de la automatización del análisis musicológico computacional}
% El análisis musicológico de grandes volúmenes de música tradicional y folclórica se ha enfrentado históricamente a una limitación metodológica: la dependencia del análisis manual de partituras, lo que supone un proceso riguroso pero difícilmente escalable, ya que la automatización digital de este procedimiento resulta prácticamente imposible.
% La digitalización de archivos musicales en formatos simbólicos estructurados (como MusicXML \citep{good2001musicxml} o MEI \citep{roland2002music}) ha abierto las puertas a la musicología computacional, pero siguen persistiendo problemas fundamentales que dificultan cualquier metodología de automatización al digitalizar este tipo de contenido.
% Estos son algunos obstáculos y dificultades encontradas en el ámbito del análisis musical que han motivado el desarrollo de este proyecto:

La motivación que impulsa esta investigación reside en la necesidad de resolver la desconexión entre el tiempo y el esfuerzo del investigador en musicología y la inabarcable escala del patrimonio digital actual. Detrás de cada corpus de música folclórica hay etnomusicólogas y documentalistas sonoros que dedican años a preservar y comprender este patrimonio cultural. Sin embargo, gran parte de su tiempo acaba absorbido por tareas repetitivas de transcripción, etiquetado y catalogación manual. Estas labores, aunque son necesarias, limitan el tiempo y la energía que podrían dedicar al análisis, la interpretación crítica y la generación de nuevo conocimiento sobre las tradiciones musicales que estudian.

El análisis etnomusicológico tradicional exige realizar análisis musicales exhaustivos y muy técnicos. Identificar la presencia de una estructura rítmica como la hemiolia, rastrear un \textit{leitmotiv} recurrente o auditar la calidad de una transcripción en un cancionero popular de miles de piezas requiere que el ojo humano recorra, compás a compás, infinidad de páginas de partituras. Este enfoque analógico es riguroso, pero carece por completo de escalabilidad. Cuando el volumen de material digitalizado alcanza el orden de las miles de obras (como ocurre en el marco de proyectos internacionales de preservación como \textit{``EA-DIGIFOLK''}, el análisis manual deja de ser una opción metodológica lenta para convertirse en una barrera física e intelectual. Ante esta realidad, muchas musicólogas se ven obligadas a trabajar con muestras reducidas que apenas representan una fracción del repertorio disponible. Esta limitación no solo restringe las preguntas que pueden plantearse, sino que también reduce la solidez estadística de sus conclusiones. Como consecuencia, una gran parte del conocimiento potencial que contienen estos archivos permanece todavía sin explorar.

A esta limitación física se suma una persistente frustración cognitiva derivada de la rigidez de los sistemas de almacenamiento actuales. La investigación científica no es un proceso lineal, sino dinámico, donde el descubrimiento de nuevos patrones y detalles genera nuevas preguntas. En el paradigma convencional, si una musicóloga formula una hipótesis a mitad de su investigación que requiera correlacionar la tesitura vocal con una temática lírica específica, se enfrenta a un muro tecnológico. Evaluar esa nueva variable implica volver a abrir, una a una, todas las partituras del corpus para extraer el dato faltante, o bien depender de un equipo técnico de desarrollo que reestructure la base de datos relacional y programe nuevas metodologías de análisis. Esta dependencia técnica rompe el flujo del pensamiento científico y somete la creatividad de la investigadora a las restricciones de los esquemas de bases de datos preexistentes.

Por tanto, la motivación de este trabajo no es sustituir el criterio de la experta por un computacional automatizado, sino liberarla de la carga operativa. Existe una necesidad urgente de diseñar herramientas que se comuniquen mediante lenguaje natural con el humanista, actuando como puentes que permitan estudiar colecciones complejas sin necesidad de escribir código ni rediseñar tablas estructuradas. Al delegar la extracción algorítmica de variables y la traducción de consultas complejas en una arquitectura de inteligencia artificial determinista, se devuelve a las musicólogas su recurso más valioso: el tiempo para pensar, interpretar y salvaguardar la riqueza cultural inmaterial de nuestros pueblos.

% El proyecto nace a partir de estas carencias, con el objetivo de resolverlas mediante el estudio y desarrollo de un sistema capaz de suplir las necesidades de las musicólogas y usuarias que necesiten herramientas más sofisticadas para enfrentarse a análisis de corpus de datos muy extensos y complejos. 

% El propósito principal de esta investigación se centra en el desarrollo de una arquitectura automatizada, robusta y escalable, capaz de analizar rasgos latentes en un corpus de obras musicales, integrando simultáneamente información estructural, rítmica, melódica y semántica para facilitar el estudio musicológico de grandes colecciones folclóricas.

\section{Objetivos}
% A partir de las limitaciones identificadas en las metodologías actuales de análisis musicológico computacional, este proyecto busca desarrollar una infraestructura tecnológica que facilite el estudio de grandes corpus musicales mediante técnicas de análisis automatizado, recuperación de información, análisis musicológico técnico y procesamiento del lenguaje natural. Para ello se establecen los siguientes objetivos.

A partir de las limitaciones críticas identificadas en las metodologías actuales de análisis musicológico (como la rigidez estructural de las bases de datos convencionales, la ineficacia de los modelos de lenguaje genéricos ante la densidad de los formatos simbólicos y la falta de escalabilidad en el análisis manual), este proyecto busca diseñar y desplegar una infraestructura tecnológica escalable y accesible para las investigadoras. Esta arquitectura híbrida tiene como objetivo unificar las técnicas de análisis musical determinista con la flexibilidad de la inteligencia artificial conversacional.

\subsection{Objetivo general}
% Desarrollar un sistema automatizado, robusto y escalable para el análisis, catalogación y consulta de grandes corpus musicales heterogéneos, capaz de extraer conocimiento musicológico a partir de un gran número de obras almacenadas en distintos formatos, facilitando así el trabajo de investigación, exploración y recuperación de información por parte de musicólogas, investigadoras y otras usuarias especializadas.

El objetivo general de este trabajo consiste en diseñar, desarrollar e implementar una arquitectura híbrida de ingeniería basada en inteligencia artificial, de ejecución local y robusta, orientada a automatizar el análisis, la catalogación sistemática y la recuperación dinámica de información en grandes corpus de música folclórica y popular. Este sistema tiene como propósito procesar y normalizar colecciones masivas de partituras digitales en formato simbólico para extraer de ellas diversas variables musicológicas complejas; estos descriptores técnicos y estructurales son posteriormente estructurados y modelados en una base de datos relacional centralizada, la cual actúa como la fuente fiable que hace determinista al sistema. El fin último de esta infraestructura es democratizar el acceso al personal no experto en este tipo de arquitecturas y facilitar el flujo de trabajo de las musicólogas e investigadoras, liberándolas de las tareas mecánicas de catalogación y de las barreras de programación tradicionales mediante una interfaz conversacional en lenguaje natural. Asimismo, con el propósito de validar la viabilidad y precisión técnica de esta propuesta frente al estado del arte, el proyecto se plantea evaluar y comparar empíricamente el rendimiento de este enfoque estructurado local frente a las capacidades de un modelo comercial de gran escala en la nube, demostrando cómo la combinación de bases de datos relacionales y protocolos de contexto específicos logra neutralizar las alucinaciones factuales y los problemas de saturación de memoria que sufren los sistemas generativos masivos ante el análisis musical especializado.

\subsection{Objetivos específicos}
Para garantizar la ejecución del propósito principal de esta investigación, se desglosan los siguientes objetivos específicos:

% \begin{itemize}
% \item Diseñar una arquitectura capaz de procesar corpus musicales compuestos por distintos formatos de representación simbólica, garantizando la compatibilidad entre las diversas fuentes documentales.
% \item Desarrollar un sistema automático de extracción y normalización de metadatos musicales con el fin de facilitar la catalogación sistemática de las obras y la construcción de bases de datos fiables y consistentes.
% \item Implementar algoritmos de análisis musical que permitan detectar y cuantificar características estructurales, melódicas, rítmicas y métricas relevantes para la investigación musicológica.
% \item Integrar el análisis musical con el análisis semántico y textual de las letras, permitiendo estudiar conjuntamente los aspectos musicales y literarios de las obras.
% \item Automatizar tareas de procesamiento y análisis que habitualmente requieren una elevada inversión de tiempo o conocimientos técnicos especializados, reduciendo las barreras de acceso a herramientas avanzadas de investigación musical.
% \item Diseñar e implementar un sistema de recuperación de información que permita realizar búsquedas complejas sobre el corpus, incluyendo consultas basadas tanto en metadatos como en características musicales extraídas automáticamente.
% \item Evaluar las capacidades y limitaciones de los modelos de lenguaje de propósito general en tareas de análisis y consulta de corpus musicológicos extensos, identificando los problemas de escalabilidad, precisión y contexto que presentan estos sistemas.
% \item Proponer mecanismos complementarios que permitan superar las limitaciones observadas en dichos modelos mediante la incorporación de análisis musicológicos especializados y estructuras optimizadas para la recuperación de conocimiento.
% \end{itemize}

\begin{itemize}
    \item \textbf{Procesamiento e ingesta multiformato.} A partir del objetivo general de admitir corpus musicales heterogéneos, se define la necesidad de diseñar e implementar un pipeline de ingesta modular capaz de parsear, validar y unificar colecciones simbólicas codificadas en los estándares \textit{MusicXML}, \textit{MEI} y \textit{MSCZ} (\textit{MuseScore}). Este subsistema debe gestional las discrepancias de codificación de las diversas fuentes documentales, garantizando una representación interna homogénea para las fases analíticas posteriores.
    \item \textbf{Extracción y normalización de metadatos estructurales.} A partir del objetivo general de automatizar la catalogación, se define el desarrollo de métodos algorítmicos encargados de aislar, normalizar y estructurar de manera sistemática los metadatos e información musical relevante de cada obra (tonalidad, compás, ámbito, procedencia geográfica, etc.). El objetivo es eliminar la inconsistencia de los registros manuales y construir una base de datos relacional sólida.
    \item \textbf{Modelado y parametrización musicológica técnica.} A partir del objetivo general de realizar un análisis musicológico exhaustivo con el fin de extraer conocimiento especializado, se define la codificación de algoritmos avanzados (explotando las capacidades analíticas y musicales de la librería \textit{music21}) dedicados a detectar, medir y cuantificar variables musicales abstractas complejas que resultan invisibles para las IAs genéricas, tales como patrones de síncopas, hemiolias, tesitura, perfiles melódicos recurrentes, etc.
    \item \textbf{Análisis híbrido músico-textual.} A partir del objetivo general de ofrecer una exploración integral del patrimonio, se define la combinación del análisis de la información musical contenida en las partituras con el estudio semántico y temático de las letras asociadas a las piezas folclóricas. Esta integración permitirá a las investigadoras relacionar características musicales específicas (como estructuras melódicas, rítmicas o armónicas) con los temas, relatos y significados presentes en la tradición oral, facilitando una interpretación más rica y contextualizada del repertorio.
    \item \textbf{Diseño de persistencia relacional optimizada.} A partir del objetivo general de garantizar un entorno determinista y escalable, se define el modelado de una arquitectura de base de datos relacional basada en \textit{SQLite}. Este motor debe estructurarse con índices óptimos y restricciones de integridad capaces de almacenar de forma eficiente todo el espectro de variables analíticas extraídas, sirviendo como la única fuente fiable del sistema.
    \item \textbf{Implementación del protocolo de comunicación local (MCP).} A partir del objetivo general de habilitar consultas flexibles sin necesidad de programar, se define el despliegue y la configuración de un servidor basado en el estándar \textit{Model Context Protocol} (MCP). Este protocolo actuará como el puente de comunicación seguro y estandarizado que conecte la base de datos relacional con la lógica del modelo de lenguaje local, dotando al sistema de la capacidad de interactuar dinámicamente con su entorno de datos.
    \item \textbf{Orquestación del motor conversacional Text-to-SQL local.} A partir del objetivo general de interactuar mediante lenguaje natural, se define el diseño de instrucciones para el modelo (\textit{prompt engineering}) y la configuración de un modelo de lenguaje de pequeña escala (\textit{Llama3.1:8B}) operado de forma local. El propósito es especializar al modelo en la traducción en tiempo real de preguntas musicológicas complejas formuladas por humanas a consultas estructuradas de SQL precisas, garantizando respuestas exactas y eliminando el riesgo de alucinación.
    \item \textbf{Estudio comparativo y evaluación empírica frente a sistemas comerciales.} A partir del objetivo general de evaluar de forma crítica las capacidades de los LLMs, se define la ejecución de un estudio empírico riguroso que contraste el rendimiento de la arquitectura local estructurada frente a un baseline comercial en la nube con RAG a través de \textit{Gemini File Search}, midiendo la precisión del modelo ante consultas de alta densidad técnica.
    \item \textbf{Desarrollo de interfaz web e interfaz de usuario.} A partir del objetivo general de democratizar el acceso a herramientas informáticas, se define el desarrollo de una aplicación web diseñada a medida que encapsule toda la complejidad técnica del sistema, las bases de datos y las conexiones MCP tras una interfaz gráfica intuitiva, eliminando cualquier fricción técnica para el usuario final.
\end{itemize}

\section{Plan de trabajo}
Aquí se describe el plan de trabajo a seguir para la consecución de los objetivos descritos en el apartado anterior.

Para la consecución de los objetivos propuestos se ha diseñado un plan de trabajo estructurado en seis fases principales, abarcando desde finales de octubre hasta principios de junio. Este cronograma contempla desde la investigación inicial y la captura de requisitos mediante trabajo de campo, hasta el desarrollo técnico de la arquitectura, su evaluación comparativa y la redacción final de la memoria de investigación.

A continuación, se detalla la planificación temporal mediante un diagrama de Gantt (Figura~\ref{fig:gantt_plan}):

\begin{figure}[htbp]
\raggedright
\hspace*{-2.5cm}
\begin{ganttchart}[
    vgrid,
    hgrid,
    x unit=0.05cm,
    y unit chart=0.7cm,
    time slot format=isodate-yearmonth,
    title/.style={fill=gray!20, draw=black, font=\small\bfseries},
    bar/.style={fill=blue!40, draw=black},
    milestone/.style={fill=red, draw=black},
    progress label text={},
    group right shift=0,
    group top shift=.6,
    group height=.2
]{2025-10}{2026-06}
    \gantttitlecalendar{year, month=letter} \\
    
    \ganttgroup{F1: Estado del Arte e Inv. Inicial}{2025-10}{2025-12} \\
    \ganttbar{Revisión bibliográfica}{2025-10}{2025-11} \\
    \ganttbar{Análisis de formatos (MusicXML, MEI)}{2025-11}{2025-12} \\
    
    \ganttgroup{F2: Diseño de la Arquitectura}{2025-12}{2026-02} \\
    \ganttbar{Diseño del esquema de base de datos SQLite}{2025-12}{2026-01} \\
    \ganttbar{Desarrollo del pipeline de normalización}{2026-01}{2026-02} \\
    
    \ganttgroup{F3: Implementación y Análisis}{2026-02}{2026-04} \\
    \ganttbar{Módulos de análisis musical (music21)}{2026-02}{2026-03} \\
    \ganttbar{Integración de análisis semántico (Ollama)}{2026-03}{2026-04} \\
    
    \ganttgroup{F4: Captura de Requisitos}{2026-03}{2026-03} \\
    \ganttbar{Entrevistas a musicólogas}{2026-03}{2026-03} \\
    \ganttbar{Definición de requisitos de usuario}{2026-03}{2026-03} \\
    
    \ganttgroup{F5: Evaluación y Pruebas}{2026-04}{2026-05} \\
    \ganttbar{Pruebas con Gemini 2.5 Flash (RAG)}{2026-04}{2026-05} \\
    \ganttbar{Evaluación comparativa Llama3 vs Gemini}{2026-05}{2026-05} \\
    
    \ganttgroup{F6: Documentación Final}{2026-04}{2026-06} \\
    \ganttbar{Redacción de la memoria técnica}{2026-04}{2026-06} \\
    % \ganttmilestone{Entrega de la Memoria}{2026-06}
\end{ganttchart}
\caption{Diagrama de Gantt del plan de trabajo (octubre - junio).}
\label{fig:gantt_plan}
\end{figure}

A nivel operativo, el desarrollo se desglosa en las siguientes actividades clave ligadas directamente a los objetivos del proyecto:
\begin{itemize}
    \item \textbf{Fase 1 (Octubre -- Diciembre):} Dedicada al estudio profundo del estado del arte en musicología computacional.
    \item \textbf{Fase 2 (Diciembre -- Febrero):} Centrada en el diseño e implementación de la infraestructura intermedia y el pipeline de normalización automática, encargados de unificar la persistencia de datos en el archivo local de SQLite independientemente de la heterogeneidad u organización original del corpus folclórico.
    \item \textbf{Fase 3 (Febrero -- Abril):} Enfocada en el desarrollo de algoritmos basados en \textit{music21} para la extracción de rasgos técnicos complejos (estructuras rítmicas y melódicas) y su integración con los modelos de lenguaje locales (\textit{Llama3} a través de \textit{Ollama}) destinados a la clasificación semántica y temática de la lírica de las obras.
    \item \textbf{Fase 4 (Marzo):} Dedicada de forma exclusiva a la captura de requisitos mediante la realización de entrevistas cualitativas a musicólogas e investigadoras del área. El objetivo de este hito es alinear las capacidades del sistema con las necesidades analíticas reales y cotidianas de las usuarias expertas.
    \item \textbf{Fase 5 (Abril -- Mayo):} Destinada a la fase crítica de evaluación. Se realizan las pruebas sobre el sistema de recuperación de información implementando tanto la aproximación local de texto a SQL como el enfoque RAG en la nube mediante la funcionalidad \textit{File Search} de la API de Gemini, permitiendo contrastar los límites de precisión, contexto y alucinaciones de ambas tecnologías.
    \item \textbf{Fase 6 (Abril -- Junio):} Consistente en la redacción de la memoria de investigación y preparación de la documentación del software desarrollado para garantizar su total reproducibilidad técnica posterior.
\end{itemize}

\section{Metodología y tecnología}
Para desarrollar este proyecto se han seleccionado herramientas de código abierto y la versión gratuita de la API de Gemini. Esta decisión garantiza la reproducción íntegra del proyecto de manera local (exceptuando la utilización de la API de Gemini).
\begin{itemize}
\item \textbf{Python 3.} Utilizado como base para desarrollar todo el proyecto. Se ha seleccionado este lenguaje de programación de alto nivel y propósito general fundamentalmente por albergar el ecosistema de librerías más avanzado para la musicología computacional, la ciencia de datos y la inteligencia artificial. Su versatilidad permite articular en un único flujo de trabajo automatizado desde el preprocesamiento de partituras y la normalización de metadatos procedentes de corpus musicales heterogéneos, hasta la gestión de bases de datos locales, la comunicación con servidores de modelos de lenguaje mediante peticiones estructuradas y la posterior evaluación analítica de los resultados.
\item \textbf{Music21.} Kit de herramientas líder para la musicología computacional. Permite realizar un análisis musical profundo de cada obra en torno a sus notas, intervalos, acordes, escalas, etc., lo que facilita el estudio de patrones armónicos, melódicos y rítmicos, difíciles de identificar con otras herramientas y, en este caso, con LLMs. Además, tiene soporte para múltiples formatos, como MIDI o XML, lo que resulta muy útil en esta investigación, donde se trabaja con un extenso y heterogéneo corpus de obras en distintos formatos. Por otra parte, su integración con Python y otras librerías facilita la extracción y el análisis de la información del corpus musical.
\item \textbf{Ollama.} En este proyecto, Ollama se despliega como un servidor local encargado de orquestar el modelo de lenguaje de pequeña escala \textit{Llama 3.1 (8B)}. Se ha escogido esta infraestructura, en primer lugar, porque garantiza el control sobre los datos y la privacidad absoluta de los corpus patrimoniales al procesar la información de manera local y, en segundo lugar, porque elimina la dependencia de costes variables por tokens y la latencia de red asociadas a servicios comerciales en la nube. Dentro de la arquitectura, este entorno se integra directamente con el protocolo avanzado \textit{Model Context Protocol} (MCP) y accede una base de datos relacional \textit{SQLite} previamente creada con información extraída de las obras del corpus. El LLM local no lee las partituras directamente, sino que actúa como un traductor algorítmico determinista (\textit{Text-to-SQL}); recibe las preguntas en lenguaje natural de las investigadoras y, gracias a las directrices de Ollama y el tipado estructurado, genera consultas SQL que se ejecutan sobre la base de datos. 
\item \textbf{API de Gemini (versión gratuita).} Plataforma de servicios en la nube basada en inteligencia artificial generativa desarrollada por Google. Representa  la contrapartida tecnológica utilizada para evaluar la validez de la propuesta local. En este proyecto se utiliza el modelo comercial \textit{Gemini 2.5 Flash} a través de su API oficial (bajo el plan de acceso para desarrollo) con el objetivo de implementar un sistema alternativo basado en RAG mediante su motor nativo \textit{File Search}. En este enfoque, los archivos del corpus musical se inyectan directamente en el índice semántico de Google, en la nube. La finalidad metodológica de la API de Gemini es estrictamente comparativa: sirve como el \textit{baseline} o estándar comercial de la industria para demostrar empíricamente cómo un modelo masivo, miles de millones de parámetros y sus inmensas ventanas de contexto, sufre de alucinaciones factuales, costes latentes y descontextualización cuando se le formulan consultas etnomusicológicas de alta complejidad técnica, en contraste con la precisión del modelo local de pequeña escala coordinado por MCP.
\item \textbf{SQLite.} Representa el motor relacional y el \textit{ground truth} determinista sobre el que se asienta la infraestructura de consulta local. A diferencia de los sistemas de gestión de bases de datos tradicionales basados en el modelo cliente-servidor (como \textit{MySQL}, \textit{PostgreSQL} o \textit{SQL Server}) (que introducen una elevada sobrecarga de administración técnico-operativa al requerir un complejo despliegue, configuraciones de red, mapeo de puertos y gestión de privilegios de acceso), \textit{SQLite} opera como un motor de base de datos relacional embebido, lo que significa que la biblioteca se integra de forma directa y nativa dentro del código de la aplicación, consolidando todo el esquema de datos, los índices y las variables musicológicas extraídas de los corpus en un único archivo físico binario local. Esta elección, en primer lugar, proporciona un rendimiento óptimo de lectura para el mapeo \textit{Text-to-SQL} gestionado por el protocolo MCP, permitiendo conexiones inmediatas y libres de latencia de red; en segundo lugar, garantiza la portabilidad absoluta del ecosistema. Al almacenar la información en un único archivo autónomo, se elimina cualquier barrera técnica de instalación para los equipos de etnomusicología, facilitando que el corpus analizado y su motor de consulta puedan ser transferidos, compartidos o respaldados con la simple copia de un directorio, cumpliendo así con el objetivo humano de democratizar y simplificar el acceso a las herramientas de computación avanzada.
\item \textbf{Lilypond.} Herramienta de notación musical de código abierto que permite generar partituras a partir de un lenguaje de marcado. Se ha utilizado para la visualización de resultados y la generación de ejemplos musicales en el proyecto.
\end{itemize}

\section{Estructura del resto del documento}
El contenido de esta memoria se organiza en diferentes capítulos de manera secuencial de acuerdo con el ciclo de vida del proyecto desarrollado:

\begin{itemize}
    \item \textbf{Capítulo 2: Estado de la Cuestión.} Presenta el marco teórico y tecnológico que sustenta la investigación. Se contextualiza el denominado \textit{giro computacional} en las humanidades digitales y la musicología, revisando los principales enfoques para el análisis automatizado de corpus musicales, la representación simbólica mediante estándares como \textit{MusicXML} y \textit{MEI}, así como las capacidades y limitaciones de los modelos de lenguaje contemporáneos aplicados al dominio musical.
    \item \textbf{Capítulo 3: Captura de requisitos y construcción del benchmark experimental.} Describe la fase de análisis de requisitos realizada mediante entrevistas y cuestionarios dirigidos a especialistas en musicología y etnomusicología. A partir de este proceso se identifican las necesidades reales de las investigadoras y se construye un conjunto de preguntas musicológicas que servirá posteriormente para evaluar las distintas arquitecturas implementadas.
    \item \textbf{Capítulo 4: Sistemas de indexación y recuperación de información basados en RAG.} Analiza las arquitecturas de recuperación aumentada por generación utilizadas como referencia experimental. Se describen tanto la implementación basada en \textit{Gemini File Search} como el sistema local construido mediante \textit{ChromaDB}, \textit{Sentence Transformers} y \textit{Ollama}. Asimismo, se estudian las limitaciones observadas en ambos enfoques cuando se enfrentan a consultas musicológicas de elevada complejidad técnica.
    \item \textbf{Capítulo 5: Diseño e implementación de la arquitectura híbrida basada en MCP y SQLite.} Constituye el núcleo metodológico de la investigación. Se presenta el pipeline completo de procesamiento del corpus, incluyendo la normalización de formatos musicales heterogéneos, la extracción automatizada de metadatos y variables musicológicas mediante \textit{music21}, el diseño de la base de datos relacional \textit{SQLite}, la implementación del servidor MCP y la integración con un modelo local (\textit{Llama 3.1:8B}) para la traducción de lenguaje natural a consultas SQL. Finalmente, se describe la aplicación web desarrollada para facilitar el acceso al sistema.
    \item \textbf{Capítulo 6: Evaluación y pruebas.} Presenta el proceso de validación de la propuesta mediante la ejecución sistemática del banco de pruebas definido durante la captura de requisitos. Se comparan los resultados obtenidos por las diferentes arquitecturas evaluadas, analizando la precisión de las respuestas, la capacidad de razonamiento musicológico y el comportamiento ante consultas complejas.
    \item \textbf{Capítulo 7: Discusión y análisis de limitaciones.} Examina los resultados desde una perspectiva cualitativa y cuantitativa. Se comparan las fortalezas y debilidades de los enfoques basados en recuperación documental frente a la arquitectura determinista sustentada en bases de datos relacionales y MCP, analizando fenómenos como las alucinaciones factuales, las limitaciones de contexto y los problemas derivados de la interpretación directa de formatos musicales simbólicos.
    \item \textbf{Capítulo 8: Conclusiones y Trabajo Futuro.} Recoge las principales aportaciones de la investigación y evalúa el grado de cumplimiento de los objetivos planteados. Asimismo, se proponen futuras líneas de desarrollo, incluyendo arquitecturas multiagente (\textit{Agent-to-Agent}), sistemas basados en \textit{GraphRAG}, modelos especializados para análisis musical y mecanismos avanzados de explotación de datos semiestructurados.
\end{itemize}

La memoria se complementa con diversos apéndices destinados a favorecer la reproducibilidad y la transferencia de conocimiento derivadas del proyecto:

\begin{itemize}

\item \textbf{Apéndice A: Actividades de transferencia y difusión científica.} Recoge las contribuciones derivadas de la investigación, incluyendo participaciones en congresos nacionales e internacionales, seminarios académicos, demostraciones técnicas y otras actividades de divulgación relacionadas con el proyecto.

\item \textbf{Apéndice B: Banco completo de pruebas y resultados experimentales.} Contiene la totalidad de las consultas utilizadas durante la evaluación empírica, así como las respuestas generadas por cada uno de los cuatro modelos analizados. Este material constituye la base experimental de la investigación y permite reproducir íntegramente el proceso de comparación descrito en los capítulos de evaluación y discusión.

\end{itemize}

Finalmente, con el objetivo de facilitar la difusión internacional de los resultados obtenidos, la memoria incorpora una versión en lengua inglesa de los capítulos correspondientes a la \textbf{Introducción} y a las \textbf{Conclusiones y trabajo futuro}, que se presentan respectivamente en los capítulos 9 y 10.

\section{Disponibilidad del código fuente y materiales de investigación}
Con el objetivo de favorecer la reproducibilidad científica y facilitar la transferencia de conocimiento derivada de esta investigación, el código fuente desarrollado durante el proyecto se encuentra disponible públicamente en un repositorio de GitHub.

El repositorio incluye la implementación completa de las distintas arquitecturas evaluadas en este trabajo, incluyendo los sistemas de recuperación aumentada por generación (\textit{RAG}), la infraestructura basada en \textit{Model Context Protocol} (MCP), los procesos de extracción automatizada de variables musicológicas mediante \textit{music21}, los esquemas de bases de datos utilizados, así como la aplicación web desarrollada para la interacción con el sistema.

El código fuente puede consultarse en el siguiente repositorio:

\begin{center}
\url{https://github.com/marsache/MIA_TFM}
\end{center}

\section{Uso de Inteligencia Artificial}
Este trabajo ha sido desarrollado con el apoyo puntual de herramientas de inteligencia artificial como asistencia técnica y de redacción. Su uso se ha limitado a tareas auxiliares, tales como la organización de ideas, la mejora de la claridad en la redacción y el apoyo en la implementación de ciertas partes del código. En todos los casos, las decisiones de diseño, la validación de resultados y la responsabilidad final sobre el contenido y la implementación recaen exclusivamente sobre la autora.


% \section{Explicaciones adicionales sobre el uso de esta plantilla}
% Si quieres cambiar el \textbf{estilo del título} de los capítulos, edita \verb|TeXiS\TeXiS_pream.tex| y comenta la línea \verb|\usepackage[Lenny]{fncychap}| para dejar el estilo básico de \LaTeX.

% Si no te gusta que no haya \textbf{espacios entre párrafos} y quieres dejar un pequeño espacio en blanco, no metas saltos de línea (\verb|\\|) al final de los párrafos. En su lugar, busca el comando  \verb|\setlength{\parskip}{0.2ex}| en \verb|TeXiS\TeXiS_pream.tex| y aumenta el valor de $0.2ex$ a, por ejemplo, $1ex$.

% TFMTeXiS se ha elaborado a partir de la plantilla de TeXiS\footnote{\url{http://gaia.fdi.ucm.es/research/texis/}}, creada por Marco Antonio y Pedro Pablo Gómez Martín para escribir su tesis doctoral. Para explicaciones más extensas y detalladas sobre cómo usar esta plantilla, recomendamos la lectura del documento \texttt{TeXiS-Manual-1.0.pdf} que acompaña a esta plantilla.

% El siguiente texto se genera con el comando \verb|\lipsum[2-20]| que viene a continuación en el fichero .tex. El único propósito es mostrar el aspecto de las páginas usando esta plantilla. Quita este comando y, si quieres, comenta o elimina el paquete \textit{lipsum} al final de \verb|TeXiS\TeXiS_pream.tex|